% Restored proof avenue: the local atlas is superseded quantitatively by
% D:uniform-local-representation-growth, but its proof is retained.
\subsection{The local minimal-extension atlas}

\begin{lemma}[Proved: a constructible atlas for reflexive modules and Frobenius]
\label{D:local-reflexive-constructible-atlas}
Let \(R,k,U\) be as in
\ref{D:uniform-local-representation-growth}, and fix a rank \(r\).
There is a finite-type parameter scheme \(T_r\) whose geometric
points give all rank-\(r\) reflexive \(R\)-modules, with a
constructible relation that is exactly ordinary isomorphism.
There is also a constructible relation that says exactly
\(M'\cong F_R^{[*]}M\).
No uniform bound on the complexity of the resulting constructible
quotient as \(r\) varies is asserted by this lemma.
\end{lemma}

\begin{proof}
Choose a projective regular resolution
\(\pi:Y\to\operatorname{Spec}(R)\) whose exceptional curves \(E_i\)
form a simple normal crossings divisor, and choose a positive
exceptional integral cycle \(Z\) with \(Z\cdot E_i<0\) for all \(i\).
Put \(I=\mathcal O_Y(-Z)\) and \(Y_n=V(I^n)\).
The \(I\)-adic and \(\mathfrak m\mathcal O_Y\)-adic topologies are
cofinal. Each \(Y_n\) is projective over \(k\).

Every reflexive \(M\) has a minimal vector-bundle extension \(F\)
on \(Y\) with the fixed slope windows
\[
 \alpha_i=g(E_i)-1\leq\mu_{\min}(F|_{E_i})
 \leq\mu_{\max}(F|_{E_i})
 \leq\beta_i=g(E_i)-1-E_i^2.
 \tag{J1}
\]
We use \cite[Definition 1.2 and Lemma 1.3, including the
elementary-transformation proof]{C_Langer2026}.
That proof applies over this complete local base: the relative
Euler characteristics used to terminate the modifications are
finite nonnegative integers, and all modifications take place
on the proper exceptional curves.
For the bundles \(F_Y^*F|_{E_i}\), the bounds of
\cite[Theorem 2.2]{C_Sun2008} and
\cite[Proposition 1.3 and its proof]{C_MehtaPauly2007} give the
enlarged windows
\[
 p\alpha_i-(p-1)\max(g(E_i)-1,0),\qquad
 p\beta_i+(p-1)\max(g(E_i)-1,0).
 \tag{J2}
\]
For genus zero the same assertion follows from splitting on
\(\mathbf P^1\). The semistable inequalities extend to arbitrary
bundles by their Harder--Narasimhan filtrations. All windows are
independent of \(r\).

\par\smallskip\noindent\textbf{Uniform poles and lifting depth.}\quad
Order the multiset of components of \(Z\).
If \(D_j\) is the prefix cycle preceding a copy of \(E_i\), the
resulting filtration of \(\mathcal O_Z\) has that layer
\(\mathcal O_{E_i}(-D_j)\).
For \(F,G\) satisfying any of the finitely many permitted windows,
\[
 H^0\bigl(E_i,\operatorname{Hom}(F_i,G_i)
                ((bZ-D_j)|_{E_i})\bigr)=0
\]
as soon as
\[
 \mu_{\max}(G_i)+(bZ-D_j)\cdot E_i<\mu_{\min}(F_i).
\]
This holds uniformly for all sufficiently large \(b\), since
\(Z\cdot E_i<0\).
The restriction sequence then shows that the global sections of
\(\operatorname{Hom}(F,G)(bZ)\) stabilize as \(b\) increases.
Every morphism on \(U\) has some finite pole order, by localization
along the Cartier divisor and quasicompactness. Thus there is one
\(B\), independent of ranks, such that every such morphism
extends as \(F\to G(BZ)\). This applies to inverse maps as well.

Choose \(N>2B\) sufficiently large that all the following \(H^1\)
groups vanish for \(n\geq N\):
\[
 H^1\bigl(Z,\operatorname{Hom}(F,G)(BZ)\otimes I^n|_Z\bigr),
 \qquad
 H^1\bigl(Z,\operatorname{End}(F)\otimes I^n|_Z\bigr).
 \tag{J3}
\]
The fixed layer filtration reduces this to smooth curves.
Serre duality identifies
\[
 H^1(E_i,\operatorname{Hom}(F_i,G_i)\otimes L)^\vee
     =\operatorname{Hom}(G_i,F_i\otimes L^{-1}\omega_{E_i}).
\]
It vanishes if
\[
 \mu_{\min}(G_i)+\deg L>
          \mu_{\max}(F_i)+2g(E_i)-2.
\]
For \(L=((B-n)Z-D_j)|_{E_i}\), its degree tends to positive
infinity with \(n\). This proves (J3), also for the Frobenius
windows (J2), without a tensor-semistability assertion.

Every vector bundle \(F_N\) on \(Y_N\) with windows (J1) lifts
successively to all \(Y_n\). The obstruction is in an \(H^2\)
group on a one-dimensional proper scheme, hence is zero.
After \(N\), the second vanishing in (J3) makes the lift unique
up to an isomorphism extending any prescribed one.
Projective Grothendieck existence algebraizes the compatible
system to a bundle \(F\) on \(Y\).
It is locally free along the exceptional fibre by faithful
flatness of completion, and everywhere by properness: a
nonempty closed nonfree locus over a local base would meet
the closed fibre. We obtain \(M=(\pi_*F)^{**}\).
This module is independent of the choice of lift up to
isomorphism. The same argument works after an algebraically
closed coefficient-field extension.
The obstruction and algebraization inputs are the
vector-bundle deformation statements and projective existence
theorem in
\cite[Tags 08MF, 08MD, 0885]{StacksLocalAtlasB}.
The first vanishing in (J3) lifts every bounded-pole map at depth
\(N\) to a map between the chosen actual lifts.

\par\smallskip\noindent\textbf{Exact isomorphism tests.}\quad
At depth \(N\), ask for maps
\[
 h_N:F_N\to G_N(BZ),\qquad k_N:G_N\to F_N(BZ)
\]
whose two products equal the canonical inclusions into the
\(2BZ\) twists.
Lift \(h_N,k_N\) by (J3).
As rational endomorphisms their products are
\(\operatorname{Id}+\varepsilon\), with
\(\varepsilon\in I^{N-2B}\operatorname{End}\).
These are invertible on the formal completion, and their
inverses algebraize by full faithfulness of existence.
Equivalently their noninvertibility loci are closed and proper
over the local base and miss its closed fibre.
Thus \(h\) is an isomorphism on \(U\), so \(M\cong M'\).
Conversely every module isomorphism gives maps with poles at most
\(BZ\), hence these equations at depth \(N\).
The finite test is therefore exact.

For Frobenius, replace \(F_N\) by \(F_{Y_N}^*F_N\).
The enlarged windows and the same lifting argument give an
equally exact test for \(M'\cong F_R^{[*]}M\).
Absolute Frobenius commutes with restriction and with base
change of a geometric point; its effect on coefficient
coordinates is their \(p\)-th power.

\par\smallskip\noindent\textbf{Finite-type parameter spaces.}\quad
Embed \(Y_N\) in a fixed projective space, with hyperplane \(H\).
The filtration of \(\mathcal O_{Y_N}\) by the component multiset
of \(NZ\) has finitely many fixed layers.
Choose \(t\) such that
\[
 \alpha_i+((t-1)H-D_j)\cdot E_i>2g(E_i)-2
\]
for each layer. Then \(H^1(Y_N,F_N(t-1))=0\).
All higher cohomology groups vanish because its support has
dimension one. Thus the ambient sheaf is \(t\)-regular and
\(F_N(t)\) is globally generated
\cite[Tag 08A8]{StacksLocalAtlasB}.
Riemann--Roch on the fixed layers gives
\(h^0(F_N(t))\leq c r\), and gives finitely many Hilbert
polynomials at fixed \(r\).

Take the finite union of Quot schemes of quotients of
\(\mathcal O_{Y_N}(-t)^a\), \(a\leq cr\), with those polynomials.
The locally free rank-\(r\) locus is open.
The conditions (J1) are open as well, using the finitely many
bad quotient Hilbert polynomials in this bounded family.
One uses all coherent quotients for their proper Quot images;
torsion only increases the degree and does not invalidate
the slope tests. This gives a finite-type scheme \(T_r\) with
a universal bundle. See \cite[Tag 0DPA]{StacksLocalAtlasB} for
bounded Quot.

On \(T_r\times T_r\), the bounded-pole pairs of maps are
parametrized by affine schemes of finite presentation, by
proper Hom representability
\cite[Tag 08K6]{StacksLocalAtlasB}.
The two inverse-product equations are closed conditions.
Their image is constructible and, by the exact lifting argument,
is precisely module isomorphism. Using the Frobenius pullback
in the source gives the other required constructible relation.
No universal algebraization of the entire parameter family
has been assumed: effectivity and equivalence were proved
on its geometric points, which is what these constructible
relations require.
\end{proof}

\begin{remark}[The independent fixed-rank route]
Applying \ref{D:effective-constructible-quotient} to this atlas
gives a constructible space of module classes and a constructible
reflexive-Frobenius function on it. Its periodic classes are
the finite-field local systems by
\ref{D:local-complete-Lang}.
Gabber's local finiteness theorem, in the precise form
\cite[Theorem 1.4]{DegliseGabber2007LocalB}, makes each finite-target
set finite. Then \ref{D:constructible-periodic-count} gives the
fixed-rank exponential bound. This argument predates and is
quantitatively superseded by the quadratic witness construction.
It is retained because it separates constructibility from the
additional rank-uniform complexity argument.
\end{remark}

